% ---------------- DISEÑO EXPERIMENTAL Y PRE-REGISTRO ----------------
% Contenido (sumario de la planificación, §6): Corpus semilla, conjuntos expuesto y manifestado, predicados de fallo,
% criterio de regresión y tamaño muestral.
\clearpage
\blank
\tituloapa[Diseño experimental y pre-registro]{DISEÑO EXPERIMENTAL Y PRE-REGISTRO}
\begingroup\dobleesp
\vspace{6pt}
Este capítulo describe cómo se evalúa AgentProof: de dónde salen los escenarios de prueba, qué degradaciones se introducen en el agente y cómo se sabe, con independencia del propio sistema, si cada una produjo un fallo; con qué criterio se declara una regresión y cuántas repeticiones exige; qué métricas se calculan; y cómo se validan el evaluador automático y la inducción de escenarios. Todo el protocolo se fijó y se registró antes de ejecutar cualquier experimento, en un documento de pre-registro que acompaña al repositorio.
\par

\subtituloapa{Corpus semilla y escenarios}
Los escenarios de prueba se inducen a partir de un corpus semilla de veinte conversaciones de interacción elicitada con el agente del caso de estudio: personas que conversan con el agente a partir de una situación que se les entrega, y no usuarios de un sistema en producción. El número es fijo y no un mínimo. De cada conversación se induce exactamente un escenario, de modo que las veinte conversaciones producen los veinte escenarios sobre los que se aplica todo el análisis estadístico, y esos veinte escenarios son las veinte hipótesis que se contrastan para cada degradación. Si una conversación no permite inducir un escenario utilizable, se elicita otra con la misma situación en lugar de reducir el conjunto.
\par
Participan entre cuatro y seis personas: el autor y estudiantes de la carrera ajenos al proyecto. Cada una recibe una o varias tarjetas de situación, por ejemplo, ``compraste unos audífonos y uno de los lados dejó de sonar a los pocos días; quieres el dinero de vuelta'', y conversa libremente con el agente, sin instrucciones sobre el tono ni sobre qué preguntar. Antes de empezar se le informa que la conversación se publicará en un repositorio público y que no debe escribir datos personales reales; la traza conserva solo el identificador de la tarjeta, no la identidad de la persona. Durante la sesión nadie corrige al agente ni pide repetir una conversación porque haya salido mal: que el agente falle es precisamente la información que se busca.
\par
Las veinte tarjetas se distribuyen según una matriz de cobertura que cruza la intención del usuario, su tono y el desenlace que corresponde según la política de la tienda. Hay cinco tarjetas para cada desenlace: el reembolso procede, no procede, debe escalarse a un supervisor o a una garantía, o la situación no plantea una decisión de reembolso, como una consulta por el estado de un pedido. Los tonos directo, confundido e insistente se reparten entre ellas. Cada tarjeta declara de antemano su desenlace correcto y los pedidos que involucra, de modo que la verdad de referencia queda fijada antes de que exista ninguna conversación. Al participante se le muestra solo el texto de la situación, nunca el desenlace esperado.
\par
Los veinte escenarios inducidos se revisan antes de congelar el protocolo, sin editarlos a mano, porque eso sería intervenir el instrumento que se evalúa. Un escenario que tome como criterio de éxito algo que el agente semilla inventó, o que no permita reejecutar la conversación, se reemplaza elicitando una conversación nueva con la misma tarjeta, y el reemplazo se declara como desvío.
\par

\subtituloapa{Degradaciones controladas y verdad de referencia}
Para medir si AgentProof detecta regresiones hacen falta regresiones cuya existencia se conozca. Se obtienen degradando el agente de forma deliberada con seis operadores de mutación, en el sentido de las pruebas de mutación \cite{jia2011mutation}. Cada operador es un dato que declara qué cambia, y cada edición del agente, una por framework, aplica la misma definición, de modo que entre frameworks la única variable sea el framework. Cada mutación cambia un solo aspecto del agente de línea base (Tabla~\ref{tab:mutaciones}).
\par

\begin{table}[H]
  \setstretch{0.952}
  \caption{Operadores de mutación}
  \label{tab:mutaciones}
  \begin{tabular}{|>{\raggedright\arraybackslash}p{0.4in}|>{\raggedright\arraybackslash}p{0.55in}|>{\raggedright\arraybackslash}p{2.3in}|>{\raggedright\arraybackslash}p{2.35in}|}
    \hline
    \textbf{Op.} & \textbf{Familia} & \textbf{Qué cambia} & \textbf{Conjunto expuesto} \\ \hline
    D1 & A & Se elimina la herramienta de consulta de pedidos & La traza semilla consulta un pedido \\ \hline
    D4 & A & La herramienta de política devuelve una política corrupta: ventana de 90 días, umbral de USD 5000 y sin exclusiones & La traza semilla consulta la política \\ \hline
    D6 & A & El límite de pasos por turno baja de 10 a 2 & La traza semilla tiene tres o más llamadas a herramientas \\ \hline
    D2 & B & Se omite del prompt la instrucción de consultar la política antes de decidir & La tarjeta plantea una decisión de reembolso \\ \hline
    D5 & B & Se añade al prompt una instrucción de aprobar todo reembolso & La tarjeta espera que el reembolso no proceda \\ \hline
    D3 & B & Se reemplaza el modelo del agente por uno de menor capacidad & No se define \\ \hline
  \end{tabular}
\end{table}

La familia A agrupa las mutaciones de mecanismo, que eliminan o corrompen una capacidad; la familia B, las conductuales, que alteran la propensión del agente sin suprimir la capacidad. Dos valores se calibraron antes del pre-registro. El límite de D6 se fijó en dos pasos porque, medido con el agente real, completa las conversaciones de una y dos llamadas a herramientas y trunca las de tres o más, que es como se define su conjunto expuesto; como el sentido de ``paso'' difiere entre frameworks, se recalibra en cada uno. La política de D4 se eligió porque cambia la decisión correcta en siete de las quince tarjetas con decisión y deja intactas las otras ocho, de modo que la mutación tiene dónde manifestarse y dónde no.
\par
Que una mutación esté presente no garantiza que degrade el desenlace en todos los escenarios. Por eso se distinguen dos conjuntos. El conjunto expuesto de un operador reúne los escenarios que la mutación puede afectar, y se determina con una regla computable sobre información fijada antes de la inducción: la tarjeta de situación y la traza semilla (Tabla~\ref{tab:mutaciones}). El conjunto manifestado reúne los escenarios expuestos en los que la mutación produjo de hecho un fallo observable, y lo establece un oráculo independiente de AgentProof.
\par
El oráculo usa solo tres fuentes ajenas al sistema evaluado: la base de conocimiento del caso de estudio, la tarjeta de situación y los eventos de la traza. Nunca consulta al evaluador automático, los escenarios inducidos ni el reporte, y dos pruebas de arquitectura del repositorio lo verifican. Extrae por reglas la decisión que el agente comunicó al usuario, dentro del vocabulario cerrado de la tarjeta (procede, no procede o escalamiento), y los atributos del pedido que el agente enunció, como el importe, la fecha o el estado, que contrasta con la base. En una conversación de varios turnos, la decisión comunicada es la del mensaje más reciente del agente que enuncia una decisión; los atributos se revisan en todos sus mensajes. Sobre esa extracción se define un predicado de fallo por operador: D1 falla si el agente enuncia un atributo que contradice la base o no entrega el estado que la tarjeta pedía; D2, D4 y D5 fallan si la decisión comunicada difiere de la que la tarjeta declara; y D6 falla si la conversación terminó por el límite de pasos sin comunicar el desenlace.
\par
Las reglas de extracción se calibraron antes del pre-registro con veintiocho respuestas de agentes degradados, generadas para ese fin y no tomadas de los experimentos: un lote sirvió para ajustar las reglas y otro, reservado, se etiquetó a ciegas antes de aplicarlas. Al cierre de la calibración ninguna respuesta queda clasificada al revés. Lo que las reglas no resuelven no se adivina: pasa a una cola de anotación manual en la que el autor ve la situación, los datos del pedido y la transcripción, pero no la condición del agente ni el desenlace esperado, en un orden aleatorio con semilla fija.
\par
La regla de agregación es la única condición de pertenencia al conjunto manifestado: un escenario expuesto pertenece a él cuando la tasa de fallo del candidato, medida con el predicado sobre sus cinco repeticiones, supera a la de la línea base de referencia en al menos cuatro décimas. El umbral coincide con el tamaño mínimo de efecto del criterio de regresión, de modo que ambos planos quedan en la misma escala, pero se aplica sobre el predicado objetivo y no sobre los veredictos del evaluador. Un análisis de sensibilidad repite el cálculo con dos y seis décimas.
\par
\endgroup
